% ═══════════════════════════════════════════════════════════════
% LERNBLATT: Repaso para la prueba oral
% Klasse 10 Spätbeginner - Zusammenfassung aller Inhalte
% ═══════════════════════════════════════════════════════════════

\documentclass[10pt, a4paper]{article}

\usepackage[utf8]{inputenc}
\usepackage[T1]{fontenc}
\usepackage[ngerman]{babel}

% --- SCHRIFTEN ---
\usepackage{helvet}
\renewcommand{\familydefault}{\sfdefault}
\usepackage{palatino}
\newcommand{\seriftext}[1]{{\fontfamily{ppl}\selectfont #1}}

\usepackage{microtype}
\usepackage[margin=1.4cm, top=1.6cm, bottom=1.3cm]{geometry}
\usepackage{enumitem}
\usepackage{tabularx}
\usepackage{booktabs}
\usepackage{array}
\usepackage{multicol}
\usepackage{tikz}

\usepackage{fancyhdr}
\usepackage{lastpage}

\usepackage{xcolor}
\usepackage{tcolorbox}
\tcbuselibrary{skins}

% ═══════════════════════════════════════════════════════════════
% FARBEN
% ═══════════════════════════════════════════════════════════════

\definecolor{spanischrot}{HTML}{C41E3A}
\definecolor{hellgrau}{HTML}{F2F2F2}
\definecolor{mittelgrau}{HTML}{777777}
\definecolor{dunkelgrau}{HTML}{444444}
\definecolor{akzent}{HTML}{C62828}

% ═══════════════════════════════════════════════════════════════
% BOXEN
% ═══════════════════════════════════════════════════════════════

\newtcolorbox{grammatikbox}{
    colback=hellgrau, colframe=hellgrau, boxrule=0pt, arc=0pt,
    left=6pt, right=6pt, top=5pt, bottom=4pt,
    borderline north={1.5pt}{0pt}{dunkelgrau}
}

\newtcolorbox{beispielbox}{
    colback=white, colframe=mittelgrau, boxrule=0.4pt, arc=0pt,
    left=6pt, right=6pt, top=4pt, bottom=4pt
}

\newtcolorbox{vokabelbox}{
    colback=white, colframe=hellgrau, boxrule=0.5pt, arc=0pt,
    left=5pt, right=5pt, top=4pt, bottom=4pt
}

% ═══════════════════════════════════════════════════════════════
% BEFEHLE
% ═══════════════════════════════════════════════════════════════

\newcommand{\es}[1]{\seriftext{\textbf{#1}}}
\newcommand{\de}[1]{{\scriptsize\textcolor{mittelgrau}{#1}}}
\newcommand{\stamm}[1]{\textcolor{akzent}{#1}}
\newcommand{\abschnitt}[1]{\vspace{6pt}\noindent\textbf{\normalsize #1}\vspace{3pt}}
\newcommand{\regel}[1]{%
    \tikz[baseline=(X.base)]\node[fill=akzent, text=white,
    font=\scriptsize\bfseries, inner sep=1.5pt, rounded corners=0.5pt](X){#1};%
}

\setlength{\parindent}{0pt}
\setlength{\parskip}{2pt}
\setlength{\columnsep}{14pt}

% ═══════════════════════════════════════════════════════════════
% KOPF- UND FUSSZEILE
% ═══════════════════════════════════════════════════════════════

\pagestyle{fancy}
\fancyhf{}
\renewcommand{\headrulewidth}{0pt}
\cfoot{\footnotesize\thepage\,/\,\pageref{LastPage}}

% ═══════════════════════════════════════════════════════════════
% DOKUMENT
% ═══════════════════════════════════════════════════════════════

\begin{document}

% --- TITEL ---
\noindent
\begin{tikzpicture}[remember picture, overlay]
    \fill[dunkelgrau] (current page.north west) rectangle +(\paperwidth, -0.9cm);
    \node[anchor=west, text=white, font=\large\bfseries]
        at ([xshift=1.4cm, yshift=-0.45cm]current page.north west) {Repaso para la prueba oral};
    \node[anchor=east, text=white, font=\small]
        at ([xshift=-1.4cm, yshift=-0.45cm]current page.north east)
        {Spanisch · Klasse 10};
\end{tikzpicture}

\vspace{0.6cm}

% ═══════════════════════════════════════════════════════════════
% DREI SPALTEN
% ═══════════════════════════════════════════════════════════════

\begin{multicols}{3}

% ---------------------------------------------------------------
% VERBOS IMPORTANTES
% ---------------------------------------------------------------
\abschnitt{1. Verbos importantes}

\begin{grammatikbox}
{\small\textbf{llamarse} \de{heißen}}\\[2pt]
{\scriptsize
\begin{tabular}{@{}ll@{}}
me llamo & nos llamamos \\
te llamas & os llamáis \\
se llama & se llaman \\
\end{tabular}}
\end{grammatikbox}

\vspace{2pt}

\begin{grammatikbox}
{\small\textbf{ser} \de{sein}}\\[2pt]
{\scriptsize
\begin{tabular}{@{}ll@{}}
soy & somos \\
eres & sois \\
es & son \\
\end{tabular}}
\end{grammatikbox}

\vspace{2pt}

\begin{grammatikbox}
{\small\textbf{tener} \de{haben}}\\[2pt]
{\scriptsize
\begin{tabular}{@{}ll@{}}
tengo & tenemos \\
\es{t\stamm{ie}nes} & tenéis \\
\es{t\stamm{ie}ne} & \es{t\stamm{ie}nen} \\
\end{tabular}}
\end{grammatikbox}

\vspace{2pt}

\begin{grammatikbox}
{\small\textbf{estar} \de{sein (Ort/Zustand)}}\\[2pt]
{\scriptsize
\begin{tabular}{@{}ll@{}}
estoy & estamos \\
estás & estáis \\
está & están \\
\end{tabular}}
\end{grammatikbox}

\vspace{4pt}

% ---------------------------------------------------------------
% VERBOS REGULARES
% ---------------------------------------------------------------
\abschnitt{2. Verbos regulares}

{\scriptsize
\begin{tabular}{@{}l@{\,}l@{\,}l@{}}
& \textbf{-ar} & \textbf{-er/-ir} \\
\midrule
yo & -o & -o \\
tú & -as & -es \\
él & -a & -e \\
nos. & -amos & -emos/-imos \\
vos. & -áis & -éis/-ís \\
ellos & -an & -en \\
\end{tabular}}

\vspace{4pt}

\begin{vokabelbox}
{\scriptsize
\textbf{-ar:} estudiar, trabajar, hablar, charlar, preparar, comprar, practicar\\[2pt]
\textbf{-er:} leer, comer, aprender, comprender\\[2pt]
\textbf{-ir:} escribir, recibir, vivir}
\end{vokabelbox}

\vspace{4pt}

% ---------------------------------------------------------------
% FRASES EJEMPLO
% ---------------------------------------------------------------
\abschnitt{3. Frases útiles}

\begin{beispielbox}
{\scriptsize
\es{Me llamo Ana y tengo 16 años.}\\[1pt]
\es{Soy de Alemania y vivo en Rostock.}\\[1pt]
\es{Estudio en el instituto.}\\[1pt]
\es{Mi padre trabaja en un hotel.}\\[1pt]
\es{Leo libros y escribo textos.}}
\end{beispielbox}

\columnbreak

% ---------------------------------------------------------------
% LA FAMILIA
% ---------------------------------------------------------------
\abschnitt{4. La familia}

{\scriptsize
\begin{tabular}{@{}ll@{}}
\es{el padre} & \es{la madre} \\
\es{el hermano} & \es{la hermana} \\
\es{el abuelo} & \es{la abuela} \\
\es{el tío} & \es{la tía} \\
\es{el primo} & \es{la prima} \\
\es{el hijo} & \es{la hija} \\
\end{tabular}}

\vspace{4pt}

\begin{beispielbox}
{\scriptsize
\es{En mi familia somos cuatro personas.}\\[1pt]
\es{Mi madre se llama María y tiene 45 años.}\\[1pt]
\es{Tengo un hermano. Se llama Pablo.}\\[1pt]
\es{No tengo hermanas.}\\[1pt]
\es{Mis abuelos viven en Berlín.}}
\end{beispielbox}

\vspace{4pt}

% ---------------------------------------------------------------
% ADJETIVOS DE PERSONALIDAD
% ---------------------------------------------------------------
\abschnitt{5. Adjetivos}

{\scriptsize
\begin{tabular}{@{}ll@{}}
\es{simpático/a} & \de{sympathisch} \\
\es{amable} & \de{freundlich} \\
\es{tranquilo/a} & \de{ruhig} \\
\es{divertido/a} & \de{lustig} \\
\es{serio/a} & \de{ernst} \\
\es{aburrido/a} & \de{langweilig} \\
\es{estricto/a} & \de{streng} \\
\es{interesante} & \de{interessant} \\
\es{inteligente} & \de{intelligent} \\
\es{tímido/a} & \de{schüchtern} \\
\es{creativo/a} & \de{kreativ} \\
\end{tabular}}

\vspace{4pt}

\begin{beispielbox}
{\scriptsize
\es{Mi padre es muy simpático.}\\[1pt]
\es{Mi hermana es un poco tímida.}\\[1pt]
\es{Soy creativo y divertido.}\\[1pt]
\es{Mi profesor es estricto pero amable.}}
\end{beispielbox}

\vspace{4pt}

% ---------------------------------------------------------------
% POSESIVOS
% ---------------------------------------------------------------
\abschnitt{6. Posesivos}

{\scriptsize
\begin{tabular}{@{}ll@{}}
\es{mi / mis} & \de{mein/e} \\
\es{tu / tus} & \de{dein/e} \\
\es{su / sus} & \de{sein/e, ihr/e} \\
\es{nuestro/a/os/as} & \de{unser/e} \\
\es{vuestro/a/os/as} & \de{euer/eure} \\
\end{tabular}}

\vspace{3pt}

{\scriptsize
\es{mi padre} · \es{mis padres} · \es{tu hermano} · \es{tus hermanos}\\[2pt]
\es{su casa} · \es{sus casas} · \es{su hijo} · \es{sus hijos}\\[2pt]
\es{nuestro padre} · \es{nuestra madre} · \es{nuestros hermanos} · \es{nuestras hermanas}\\[2pt]
\es{vuestro profesor} · \es{vuestra profesora} · \es{vuestros amigos} · \es{vuestras amigas}}

\columnbreak

% ---------------------------------------------------------------
% EL BARRIO
% ---------------------------------------------------------------
\abschnitt{7. El barrio}

{\scriptsize
\begin{tabular}{@{}ll@{}}
\es{la biblioteca} & \es{la piscina} \\
\es{la panadería} & \es{el supermercado} \\
\es{el parque} & \es{la plaza} \\
\es{la calle} & \es{el mercado} \\
\es{la farmacia} & \es{la peluquería} \\
\es{el polideportivo} & \\
\end{tabular}}

\vspace{3pt}

{\scriptsize Adj.: \es{tranquilo}, \es{ruidoso} \de{laut}, \es{antiguo}, \es{nuevo}, \es{grande}, \es{pequeño}}

\vspace{4pt}

% ---------------------------------------------------------------
% HAY + ESTAR
% ---------------------------------------------------------------
\abschnitt{8. Hay y estar}

\begin{grammatikbox}
{\scriptsize
\regel{hay} \de{es gibt}\\[3pt]
\es{hay} + art. indef.: \es{Hay \underline{una} biblioteca.}\\
\es{hay} + número: \es{Hay \underline{tres} parques.}\\
\es{hay} + cantidad: \es{Hay \underline{muchas} tiendas.}\\
\es{hay} + sin art.: \es{Hay cafés y restaurantes.}}
\end{grammatikbox}

\vspace{3pt}

\begin{grammatikbox}
{\scriptsize
\regel{estar} \de{wo?}\\[3pt]
\es{La biblioteca \underline{está} en la calle Mayor.}\\
\es{El parque \underline{está} cerca de mi casa.}\\
\es{¿Dónde \underline{está} el supermercado?}}
\end{grammatikbox}

\vspace{4pt}

\begin{beispielbox}
{\scriptsize
\es{En mi barrio hay un parque grande.}\\[1pt]
\es{Hay muchas tiendas y cafés.}\\[1pt]
\es{La piscina está cerca del instituto.}\\[1pt]
\es{Mi barrio es tranquilo y bonito.}}
\end{beispielbox}

\vspace{4pt}

% ---------------------------------------------------------------
% PREGUNTAS
% ---------------------------------------------------------------
\abschnitt{9. Preguntas}

{\scriptsize
\begin{tabular}{@{}ll@{}}
\es{¿Qué?} & \de{Was?} \\
\es{¿Quién?} & \de{Wer?} \\
\es{¿Cómo?} & \de{Wie?} \\
\es{¿Dónde?} & \de{Wo?} \\
\es{¿De dónde?} & \de{Woher?} \\
\es{¿Cuál?} & \de{Welche/r?} \\
\es{¿Cuánto/a/os/as?} & \de{Wie viel/e?} \\
\end{tabular}}

\vspace{3pt}

{\scriptsize
\es{¿Cómo te llamas?} · \es{¿De dónde eres?}\\
\es{¿Cuántos años tienes?} · \es{¿Dónde vives?}\\
\es{¿Cómo es tu madre?} · \es{¿Qué hay en tu barrio?}}

\vspace{4pt}

% ---------------------------------------------------------------
% PREPOSICIONES
% ---------------------------------------------------------------
\abschnitt{10. Preposiciones}

{\scriptsize
\es{de} \de{aus, von}: \es{Soy \underline{de} Alemania.}\\
\es{en} \de{in}: \es{Vivo \underline{en} Rostock.}\\
\es{con} \de{mit}: \es{Vivo \underline{con} mi familia.}}

\end{multicols}

% ═══════════════════════════════════════════════════════════════
% SEITE 2: FRASES COMPLETAS
% ═══════════════════════════════════════════════════════════════

\newpage

\begin{center}
{\large\bfseries Frases completas para la prueba oral}\\[3pt]
{\small Vollständige Sätze für die mündliche Prüfung}
\end{center}

\vspace{6pt}

\begin{multicols}{2}

% ---------------------------------------------------------------
% PRESENTARSE
% ---------------------------------------------------------------
\abschnitt{Presentarse} \de{Sich vorstellen}

\begin{beispielbox}
{\small
\es{¡Hola! Me llamo [Name] y tengo [X] años.}\\[2pt]
\es{Soy de Alemania y vivo en [Stadt].}\\[2pt]
\es{Soy alumno/alumna en el instituto [Name].}\\[2pt]
\es{Estudio español, inglés, matemáticas...}\\[2pt]
\es{Soy [simpático/a, divertido/a, creativo/a...].}}
\end{beispielbox}

\vspace{6pt}

% ---------------------------------------------------------------
% LA FAMILIA
% ---------------------------------------------------------------
\abschnitt{Hablar de la familia} \de{Über die Familie sprechen}

\begin{beispielbox}
{\small
\es{En mi familia somos [X] personas.}\\[2pt]
\es{Tengo [un hermano / una hermana / dos hermanos].}\\[2pt]
\es{No tengo hermanos. Soy hijo/a único/a.}\\[2pt]
\es{Mi padre se llama [Name] y tiene [X] años.}\\[2pt]
\es{Mi madre trabaja en [un hospital / una oficina...].}\\[2pt]
\es{Mis abuelos se llaman [Name] y [Name].}\\[2pt]
\es{Mi abuela tiene [X] años y es muy simpática.}\\[2pt]
\es{Tengo [dos primos / una prima / un tío...].}}
\end{beispielbox}

\vspace{6pt}

% ---------------------------------------------------------------
% DESCRIBIR PERSONAS
% ---------------------------------------------------------------
\abschnitt{Describir personas} \de{Personen beschreiben}

\begin{beispielbox}
{\small
\es{Mi padre es alto y tiene el pelo negro.}\\[2pt]
\es{Mi madre es muy amable y simpática.}\\[2pt]
\es{Mi hermano es un poco tímido pero muy inteligente.}\\[2pt]
\es{Mi mejor amigo/a es divertido/a y creativo/a.}\\[2pt]
\es{Mi profesor/a de español es estricto/a pero interesante.}}
\end{beispielbox}

\columnbreak

% ---------------------------------------------------------------
% EL BARRIO
% ---------------------------------------------------------------
\abschnitt{Describir el barrio} \de{Das Stadtviertel beschreiben}

\begin{beispielbox}
{\small
\es{Vivo en un barrio [tranquilo / grande / pequeño].}\\[2pt]
\es{En mi barrio hay [un parque, una biblioteca, tiendas...].}\\[2pt]
\es{Hay muchos cafés y restaurantes.}\\[2pt]
\es{No hay piscina, pero hay un polideportivo.}\\[2pt]
\es{La biblioteca está en la calle [Name].}\\[2pt]
\es{El supermercado está cerca de mi casa.}\\[2pt]
\es{Mi barrio es bonito y tiene muchas zonas verdes.}}
\end{beispielbox}

\vspace{6pt}

% ---------------------------------------------------------------
% QUÉ SE PUEDE HACER
% ---------------------------------------------------------------
\abschnitt{Qué se puede hacer} \de{Was man machen kann}

\begin{beispielbox}
{\small
\es{En mi barrio se puede pasear por el parque.}\\[2pt]
\es{Puedo comprar en el supermercado.}\\[2pt]
\es{Hay una biblioteca donde puedo leer libros.}\\[2pt]
\es{En el polideportivo practico deporte.}\\[2pt]
\es{Me gusta pasear con mis amigos en la plaza.}}
\end{beispielbox}

\vspace{6pt}

% ---------------------------------------------------------------
% RESPONDER A PREGUNTAS
% ---------------------------------------------------------------
\abschnitt{Responder a preguntas}

\begin{beispielbox}
{\small
\es{-- ¿Estudias o trabajas?}\\
\es{-- Estudio en el instituto.}\\[4pt]
\es{-- ¿Vives en una ciudad o en un pueblo?}\\
\es{-- Vivo en una ciudad / en un pueblo pequeño.}\\[4pt]
\es{-- ¿Qué hay en tu barrio?}\\
\es{-- Hay un parque, tiendas y una biblioteca.}\\[4pt]
\es{-- ¿Dónde está el supermercado?}\\
\es{-- Está en la calle Mayor, cerca del parque.}}
\end{beispielbox}

\end{multicols}

\vspace{8pt}

% ═══════════════════════════════════════════════════════════════
% FUSSZEILE
% ═══════════════════════════════════════════════════════════════

\begin{center}
\rule{0.8\textwidth}{0.4pt}\\[4pt]
{\small\textit{¡Practica en voz alta! -- Übe laut sprechen!}}
\end{center}

\end{document}
